\documentclass[a4paper,12pt, twoside]{article}
\usepackage[a4paper, left=2.5cm, right=2.5cm, top=2.5cm, bottom=2.5cm, bindingoffset=0.5cm]{geometry}
\usepackage[utf8x]{inputenc}
\usepackage{polski}
\usepackage[polish]{babel}
\usepackage{graphicx}
\graphicspath{ {C:/Users/kuban/Desktop/Inzynierka/E-Card/Praca dyplomowa/images/} }
\usepackage{indentfirst}
\usepackage{float}
\usepackage{caption}
\usepackage{adjustbox}
\usepackage{array}
\usepackage{makecell}
\usepackage{amsmath}
\usepackage{enumitem}
\usepackage{afterpage}
\usepackage{subfig}
\usepackage{hyperref}
\usepackage{url}
\renewcommand\thesection{\arabic{section}.}
\renewcommand\thesubsection{\arabic{section}.\arabic{subsection}.}
\renewcommand\thesubsubsection{\arabic{section}.\arabic{subsection}.\arabic{subsubsection}.}
\frenchspacing
\setlength{\parindent}{.5cm}
\makeatletter
\setlength{\@fptop}{0pt}
\makeatother
\linespread{1.5}
\begin{document}
	\newpage
	\thispagestyle{empty}
	\begin{center}
		
		\begin{figure}
			\centering
			\includegraphics[width=5cm]{logo_polsl.jpg}
			\vspace{.5cm}
		\end{figure}
		
		{\fontsize{17}{17}\selectfont
			\textsc{Politechnika Śląska \\[.3cm]
				Wydział Automatyki, Elektroniki i Informatyki  \\[.3cm]
				Kierunek Automatyka i Robotyka  \\[1.5cm]}
			\textbf{Projekt inżynierski \\[0.7cm]}}
		
		\Large
		{Elektroniczna plakietka konfigurowana z urządzenia mobilnego \\[4cm]}
		\Large{\begin{flushleft}
				Autor: Jakub Legutko\\
				Kierujący pracą: dr inż. Grzegorz Dziwoki\\[0.3cm]
		\end{flushleft}}
		
		\normalsize
		\vfill Gliwice, styczeń 2019
	\end{center}
	\newpage
	\newpage
	\thispagestyle{empty}
	\tableofcontents
	\newpage
	%\leavevmode\thispagestyle{empty}\newpage
	\newpage
	\clearpage
	\setcounter{page}{1}
	\section{Wstęp}
	
	\subsection{Wprowadzenie}
	Zmiany klimatyczne są coraz bardziej widoczne na świecie. Wycinka drzew jest jednym z głównych powodów wzrostu emisji gazów cieplarnianych\cite{clima_causes}, ograniczenie jej wpłynie pozytywnie na poziom ${CO_{2}}$ w atmosferze, co przełoży się na polepszenie klimatu. 
	
	Druk wizytówek oraz identyfikatorów pochłania znaczące ilości papieru, oraz tuszu w skali globalnej ze względu na ich jednorazowe wykorzystanie w większości przypadków. Z tego powodu ważne jest znalezienie alternatywnej metody. Gdyby zamiast papieru wykorzystać urządzenie, z którego można korzystać wielokrotnie poprzez zmianę wyświetlanych danych, zmniejszylibyśmy zapotrzebowanie na papier, zmniejszając jednocześnie emisję dwutlenku węgla. Oprócz papieru potrzebnego do wytworzenia plakietek nie będziemy również marnować toneru, którego produkcja pochłania około 3.2 kilograma gazów cieplarnianych\cite{cartidge_production} na kartridż zawierający 200 gramów toneru. Kolejną rzeczą jest recykling, któremu trzeba poddać niepotrzebne już identyfikatory oraz zużyte do produkcji ich kartridże.
	
	Idealnym rozwiązaniem takiego problemu wydaje się papier elektroniczny (ang. \textit{e-paper}). Wyświetlane na nim dane można modyfikować dowolną ilość razy, dzięki czemu może być wykorzystywany wielokrotnie, nie powodując zużycia dodatkowych surowców. Po wyświetleniu dostarczonych danych nie jest już zużywana energia. Może przez to być używany bez dodatkowych baterii lub innego dodatkowego źródła zasilania.
	\newpage
	
	
	\subsection{Cel projektu}
	
	
	
	\section{Podsumowanie}
	
	\newpage
	\section{Bibliografia}
	
	\begingroup
	\renewcommand{\section}[2]{}%
	\begin{thebibliography}{}
		
		\bibitem{clima_causes}
		https://ec.europa.eu/clima/change/causes\_pl\\
		Dostęp: 23 grudnia 2019
		
		\bibitem{cartidge_production}
		https://www.energycentral.com/c/ec/ink-waste-environmental-impact-printer-cartridges\\
		Dostęp: 24 grudnia 2019
		
		
	\end{thebibliography}
	\endgroup
	
	\newpage
	\section{Spis rysunków}
	\begingroup
	\renewcommand{\section}[2]{}%
	\listoffigures
	\endgroup
	
\end{document}

